\section{条件互信息与 Markov 结构}
\label{sec:07-Information-Theory/02-Joint-Information/03}

你早上出门，看见街上有人撑伞。你猜路是湿的。这个推理有效，因为撑伞和路湿都跟下雨有关。可如果我已经站在窗边看见了倾盆大雨，再转头看街上有没有人撑伞，对判断路湿还有帮助吗？

信息论对这类情景有一个量化工具：把互信息放进条件层，问已知某个变量后，另两个变量之间还剩多少统计关联。答案不总是``变少''。有些结构里，条件化消灭依赖；另一些结构里，条件化反而制造依赖。

\subsection{一个数字例子：先不算条件，只算依赖}

把刚擦的情景换成数字。设天气 \(Z\) 有两种状态：下雨（\(z_1\)）概率 \(0.4\)，晴天（\(z_2\)）概率 \(0.6\)。撑伞 \(X\) 和路湿 \(Y\) 都概率化地依赖天气：

\[
\begin{aligned}
P(X=\text{撑}\mid z_1) &= 0.9, & P(X=\text{撑}\mid z_2) &= 0.1,\\
P(Y=\text{湿}\mid z_1) &= 0.95, & P(Y=\text{湿}\mid z_2) &= 0.05.
\end{aligned}
\]

不仅如此，我们假设给定天气后，撑伞和路湿各自独立触发——下雨既让人撑伞，也让路变湿，但撑伞的人不会额外弄湿路面。有了这些数字，可以先算边缘分布，再算 \(I(X;Y)\)。\(X\) 和 \(Y\) 在总体上不是独立的：雨天两项都高，晴天两项都低，所以互信息为正。具体数字取决于你代入公式后算出的结果，但方向是确定的——观察到撑伞，你对路湿的判断会改变。

这就是日常推理的信息基础。但问题卡在这里：这个依赖是否只是天气的投影？

\subsection{换个问法：知道天气之后呢}

现在设想你已经知道了 \(Z=z_1\)（下雨）。在这个条件世界里，撑伞概率是 \(0.9\)，路湿概率是 \(0.95\)，而且我们假设两者条件独立。那么给定 \(Z=z_1\) 后，\(X\) 和 \(Y\) 的联合分布就是乘积：

\[
P(x,y\mid z_1) = P(x\mid z_1)\,P(y\mid z_1).
\]

互信息在条件分布上的值应当为零。同样，在 \(Z=z_2\) 的世界里也是零。平均起来，已知天气后，撑伞不再提供关于路湿的额外线索。

条件互信息就是把互信息放进每个 \(Z=z\) 的切片里再平均：

\[
I(X;Y\mid Z)
= \sum_z P(z)\, I(X;Y\mid Z=z)
= H(X\mid Z) - H(X\mid Y,Z).
\]

它衡量已知 \(Z\) 后，\(Y\) 还能为 \(X\) 消除多少不确定性。在我们的例子里，这个量是零。撑伞之所以让你猜到路湿，是因为它透露了天气；天气已知，这条间接信息就不再新增任何东西。

\subsection{KL 形式让条件独立变成一个等式}

把互信息写成 KL 散度的期望，条件版本同样紧凑：

\[
I(X;Y\mid Z) = \E_Z\Bigl[
D\bigl(p_{XY\mid Z} \,\|\, p_{X\mid Z}\,p_{Y\mid Z}\bigr)
\Bigr].
\]

内层是条件分布与``条件独立乘积''之间的 KL 散度。KL 非负，期望保持非负，所以

\[
I(X;Y\mid Z) \ge 0,
\]

等号成立当且仅当对几乎所有 \(z\)，有 \(p_{XY\mid Z} = p_{X\mid Z}\,p_{Y\mid Z}\)——即

\[
X \perp Y \mid Z.
\]

概率论里的条件独立，在信息论里变成了一个精确的零信息量陈述。这比单独用概率不等式检验条件独立更统一：互信息的链式法则和上下界都可以直接套用。

\subsection{共同原因：条件化让互信息消失}

天气同时影响撑伞和路湿。这种结构叫共同原因：\(X \leftarrow Z \rightarrow Y\)。总体数据里 \(X\) 和 \(Y\) 正相关，但给定 \(Z\) 后相关消失。

这不是特例。任何满足 \(X \leftarrow Z \rightarrow Y\) 且给定 \(Z\) 后 \(X,Y\) 独立的生成过程，都有

\[
I(X;Y) > 0, \qquad I(X;Y\mid Z) = 0.
\]

统计里把它叫作可解释的混杂：\(X\) 和 \(Y\) 的关联完全由 \(Z\) 传导。控制 \(Z\) 后，剩余关联为零。

\subsection{共同结果：条件化反而制造依赖}

现在把箭头反过来。设 \(X\) 和 \(Y\) 是两枚独立公平硬币，各以 \(1/2\) 概率取 \(0\) 或 \(1\)。定义

\[
Z = X \oplus Y,
\]

即异或结果。无条件时，\(X\) 和 \(Y\) 独立：

\[
I(X;Y) = 0.
\]

但现在给定 \(Z=1\)。如果 \(X=0\)，则 \(Y\) 必须是 \(1\)；如果 \(X=1\)，则 \(Y\) 必须是 \(0\)。在 \(Z=1\) 的条件世界里，\(X\) 完全确定 \(Y\)。类似地，给定 \(Z=0\) 时也一样。因此在每个条件切片里互信息都是一 bit，平均后：

\[
I(X;Y\mid Z) = 1 \text{ bit}.
\]

条件化从零互信息出发，凭空制造了一 bit 依赖。

直觉上：两个独立原因产生同一个结果，一旦你知道了结果，原因之间就被``锁死''了。这就是统计里的对撞偏差——控制共同结果会打开本来不存在的关联路径。

\subsection{不能随便条件化}

上述两个例子合在一起揭示了一条操作原则：加入或删除条件变量，可能消灭依赖，也可能制造虚假依赖。共同原因结构里控制 \(Z\) 消除混杂；共同结果结构里控制 \(Z\) 反而引入选择偏差。

机械地``把所有能测的变量都放进条件集''不是安全策略。条件互信息提供了量化工具——你可以算 \(I(X;Y\mid Z)\) 来判断条件化后的剩余关联——但因果方向仍需要结构假设和领域知识。信息论告诉你关联的强弱，不替你决定哪些关联反映真实的传导关系。

\subsection{互信息链式法则：谁先提供、谁后补充}

回到信息量本身。\((Y,Z)\) 联合起来关于 \(X\) 的信息，可以拆成两步：

\[
I(X;Y,Z) = I(X;Z) + I(X;Y\mid Z).
\]

证明只需展开定义：

\[
\begin{aligned}
I(X;Z) + I(X;Y\mid Z)
&= [H(X) - H(X\mid Z)] + [H(X\mid Z) - H(X\mid Y,Z)] \\
&= H(X) - H(X\mid Y,Z) = I(X;Y,Z).
\end{aligned}
\]

中间项 \(H(X\mid Z)\) 一加一减，望远镜相消。这条链式法则把总信息分成两段：\(Z\) 先提供的那部分，以及已知 \(Z\) 后 \(Y\) 额外贡献的那部分。

顺序任意。你也可以先让 \(Y\) 出场：

\[
I(X;Y,Z) = I(X;Y) + I(X;Z\mid Y).
\]

两个分解等号成立，但不是同一种拆法。选择哪个顺序，取决于你实际获得信息的先后——先做血液检查再看影像，与先拍片再抽血，信息增量可以相同，但中间每一步的``已知条件''不同。

推广到序列：

\[
I(X;Y_1,\ldots,Y_n) = \sum_{i=1}^{n} I(X;Y_i \mid Y_1,\ldots,Y_{i-1}).
\]

这条公式是逐步观测、主动学习和自适应通信的信息账本：每加入一个新信号，链式法则告诉你它在已有知识之上净增加了多少关于 \(X\) 的线索。如果一个信号的条件互信息已经接近零，说明前面的观测已经把相关信息榨干了。

\subsection{Markov 链：信息阻断的精确条件}

设三个变量排成顺序

\[
X \to Y \to Z,
\]

含义是给定 \(Y\) 后，\(Z\) 与 \(X\) 条件独立：

\[
p(z \mid x, y) = p(z \mid y).
\]

换句话说，\(Z\) 不直接依赖 \(X\)，它只通过 \(Y\) 间接联系 \(X\)。信息论等价表述是

\[
I(X;Z \mid Y) = 0.
\]

Markov 结构的联合分布具有简洁的因式分解：

\[
p(x,y,z) = p(x)\,p(y\mid x)\,p(z\mid y),
\]

不出现 \(p(z\mid x)\) 这样的直接项。注意这里的箭头是条件独立关系的图形标记，不自动等于因果方向。在通信流水线或神经网络前馈中，它常有自然的因果解释；但仅从联合分布反推，同一个分布可能兼容多种箭头方向。

Markov 条件对应信息流的一个硬性阻断：\(Y\) 一旦已知，\(X\) 的过去就不再影响 \(Z\) 的未来。这个性质是下一节数据处理不等式的地基。

\subsection{小结：条件化如何重塑信息版图}

条件互信息告诉你，在已知某些变量之后，剩余的相关性还剩多少。它可以为零（依赖被解释）、可以为正（仍有直接关联），也可以从零变正（条件化制造关联）。Markov 结构和链式法则为这些判断提供了运算骨架。

但截至目前，我们只说了条件化``改变''互信息。一个更具体的操作问题在等着：如果你主动处理数据——压缩、变换、加噪——关于原始信号的信息会发生什么？下一节把 Markov 链放进处理流水线，给出一个看似简单但后果深远的不等式。
